% ============================================================
%  Capítulo 6 — Resultados
% ============================================================

\section{Resultados}
\label{sec:resultados}

Este capítulo valida empíricamente las cuatro afirmaciones centrales de la
arquitectura propuesta: (i)~el controlador determinista orquesta el flujo de forma
robusta; (ii)~el agente planificador selecciona modelos de forma adaptativa a partir
del \textit{pool} de experiencias; (iii)~el \textit{pool} acumula priors empíricos
que mejoran decisiones futuras; y (iv)~las puertas de supervisión humana son
suficientemente ligeras porque los agentes producen planes de alta calidad.

Si bien el sistema admite tres tipos de problema —clasificación, regresión y
pronóstico de series temporales—, el núcleo técnico de esta tesis es el pronóstico.
Es en este dominio donde la selección adaptativa de modelos resulta más valiosa: la
diversidad de regímenes (estadístico, supervisado con variables exógenas,
paseo~aleatorio) hace que ningún modelo único sea óptimo para todos los casos, y donde el
\textit{pool} aporta el mayor valor diferencial respecto a un enfoque sin memoria.
La clasificación y la regresión se incluyen para demostrar que la arquitectura es
agnóstica al dominio, pero su tratamiento en este capítulo es secundario.

La estructura del capítulo es la siguiente.
La Sección~\ref{sec:pool-evaluacion} evalúa la calidad del \textit{pool} de
experiencias a partir del conjunto de benchmarks ejecutado antes de las pruebas con
usuario.
La Sección~\ref{sec:casos-estudio} presenta seis casos de estudio a través del
pipeline completo en modo interactivo, cada uno diseñado para validar un aspecto
específico del comportamiento agéntico.
Finalmente, la Sección~\ref{sec:coste-agente} cuantifica el coste temporal y de
tokens de los componentes agénticos frente a los deterministas.

% ============================================================
\subsection{Evaluación del \textit{pool} de experiencias}
\label{sec:pool-evaluacion}
% ============================================================

Antes de ejecutar cualquier experimento interactivo, el sistema se inicializó con 23
registros de experiencia obtenidos mediante el ejecutor de benchmarks
(\texttt{run\_benchmark.py}), que procesa cada \textit{dataset} con la misma pila
determinista que utiliza el pipeline de producción.
En el momento de escritura de esta tesis, el \textit{pool} contiene: 11 registros de
pronóstico, 7 de clasificación y 5 de regresión.
Los benchmarks se ejecutaron con 4 ensayos de optimización de Optuna por modelo,
usando validación cruzada en ventana expandida para pronóstico y validación cruzada
estratificada de 5 particiones para clasificación y regresión.

Las subsecciones siguientes presentan los resultados por tipo de problema.

% ── 6.1.1  Pronóstico ──────────────────────────────────────────────────────

\subsubsection{Pronóstico de series temporales}
\label{sec:pool-forecasting}

La Tabla~\ref{tab:pool-forecasting} recoge los 11 \textit{datasets} de pronóstico
del \textit{pool}.
Para cada uno se indica la frecuencia temporal, el horizonte de predicción, el número
de observaciones, el modelo seleccionado como campeón por el ejecutor de
entrenamiento, el error cuadrático medio sobre el conjunto de validación (RMSE) y el
error normalizado nRMSE = $\text{RMSE} / \bar{y} \times 100$, que permite comparar
la magnitud del error de forma adimensional entre series de escalas muy distintas.
La columna \textit{Régimen} indica si el \textit{dataset} tiene una etiqueta de
familia esperada en el manifiesto de benchmarks, y la columna \textit{Conc.}\ señala
si el modelo campeón pertenece a dicha familia.

\begin{table}[H]
\centering
\caption{Resultados del benchmark de pronóstico (11 \textit{datasets}).
 ETS = suavizado exponencial automático; ARIMA = AutoARIMA;
 SVR = máquinas de vectores de soporte para regresión;
 BA = bosque aleatorio; XGB = XGBoost; ET = \textit{Extra Trees}; GBM =
 \textit{Gradient Boosting}.
 nRMSE = $\text{RMSE}/\bar{y}\times 100$.
 {\checkmark} = familia campeón concuerda con la etiqueta esperada;
 {\ding{55}} = discrepancia; {---} = sin etiqueta esperada.}
\label{tab:pool-forecasting}
\small
\setlength{\tabcolsep}{4pt}
\begin{tabular}{llrrrrrrll}
\toprule
\textbf{Dataset} & \textbf{Frec.} & \textbf{H} & \textbf{Filas}
  & \textbf{Modelo campeón} & \textbf{Familia}
  & \textbf{RMSE} & \textbf{nRMSE\,(\%)}
  & \textbf{Régimen} & \textbf{Conc.} \\
\midrule
\multicolumn{10}{l}{\textit{Régimen estadístico (etiquetados)}} \\[2pt]
Air Passengers         & MS  & 12 &    144 & AutoARIMA       & Estadístico &   23,56 &  8,4 & Estadístico & \checkmark \\
CO$_2$ Mauna Loa       & MS  &  6 &    526 & ETS             & Estadístico &    0,33 &  0,1 & Estadístico & \checkmark \\
Manchas solares        & YS  &  3 &    309 & SVR             & ML~sup.     &    9,55 & 19,2 & Estadístico & \ding{55}  \\
Río Nilo               & YS  &  3 &    100 & SVR             & ML~sup.     &  114,53 & 12,5 & Estadístico & \ding{55}  \\
\midrule
\multicolumn{10}{l}{\textit{Régimen supervisado (etiquetado)}} \\[2pt]
Metro Minneapolis      &  D  & 14 &  1.860 & Bosque Aleatorio & ML~sup.     &  240,72 &  7,4 & Supervisado & \checkmark \\
\midrule
\multicolumn{10}{l}{\textit{Sin etiqueta esperada}} \\[2pt]
Electricidad Victoria  &  D  & 14 &  4.807 & ETS             & Estadístico & 27.488  &  9,9 & ---         & ---        \\
S\&P~500               &  W  & 13 &  1.011 & XGBoost         & ML~sup.     &  177,98 &  7,9 & ---         & ---        \\
Petróleo WTI           &  W  & 13 &  1.013 & ETS             & Estadístico &    5,56 &  7,7 & ---         & ---        \\
Oro + Macro            &  W  & 13 &  1.013 & Bosque Aleatorio & ML~sup.     &   65,26 &  5,0 & ---         & ---        \\
Bitcoin (BTC)          &  W  & 13 &    335 & Extra Trees     & ML~sup.     & 4.483   & 18,5 & ---         & ---        \\
EUR/USD                &  W  & 13 &    335 & GBM             & ML~sup.     &    0,02 &  1,6 & ---         & ---        \\
\bottomrule
\end{tabular}
\end{table}

\paragraph{Análisis por régimen.}
De los cinco \textit{datasets} con etiqueta de familia esperada, tres coinciden con la
predicción (60\,\%), y dos —Manchas Solares y Río Nilo— presentan discrepancia: el
modelo campeón elegido en validación es SVR en lugar de un método estadístico clásico.
Este resultado no es un fallo del sistema sino un hallazgo legítimo: ambas series son
cortas (100–309 observaciones), carecen de variables exógenas y presentan
irregularidades no capturables por ETS ni AutoARIMA en el horizonte de validación;
el SVR, al no asumir ninguna estructura paramétrica, consigue una generalización
marginalmente superior en estas condiciones.
El \textit{pool} almacena el campeón real, no el esperado; las etiquetas de régimen
son únicamente un prior para el planificador, no una restricción.

Los cinco \textit{datasets} financieros (S\&P~500, Petróleo, Oro, Bitcoin,
EUR/USD) no tienen etiqueta esperada porque su comportamiento es ambiguo en el
histórico: la hipótesis del mercado eficiente predice que modelos naive deberían ganar,
pero en el backtesting con ventana histórica de hasta 19 años, los modelos de árboles
y ETS logran capturar autocorrelaciones residuales.
El \textit{pool} registra fielmente estos campeones sin forzar ninguna clasificación
a priori.

Todos los valores nRMSE se encuentran por debajo del 20\,\%, salvo Manchas Solares
(19,2\,\%) y Bitcoin (18,5\,\%), que son series de gran variabilidad intrínseca.
El nRMSE más bajo corresponde a CO$_2$ Mauna Loa (0,1\,\%), cuya tendencia lineal
+ ciclo anual es casi perfectamente capturable por ETS.

% ── 6.1.2  Clasificación y regresión ───────────────────────────────────────

\subsubsection{Clasificación y regresión}
\label{sec:pool-clasif-regr}

Las Tablas~\ref{tab:pool-classification} y~\ref{tab:pool-regression} recogen los
resultados del benchmark para los tipos de problema secundarios.
Su inclusión en el \textit{pool} garantiza que el sistema responda de forma coherente
cuando el usuario sube un \textit{dataset} tabular no temporal, demostrando que la
arquitectura agente–planificador–ejecutor es agnóstica al dominio.

\begin{table}[H]
\centering
\caption{Resultados del benchmark de clasificación (7 \textit{datasets}).
 Métrica: macro-F1 sobre el conjunto de validación (5-fold).
 RL = regresión logística; BA = bosque aleatorio; LGBM = LightGBM.}
\label{tab:pool-classification}
\small
\begin{tabular}{llrll}
\toprule
\textbf{Dataset} & \textbf{Dominio} & \textbf{Filas} & \textbf{Modelo campeón} & \textbf{macro-F1} \\
\midrule
Iris                  & Botánica           &     150 & Regresión Logística & 0,966 \\
Vino                  & Enología           &     178 & Bosque Aleatorio    & 0,980 \\
Cáncer de mama        & Medicina           &     569 & LightGBM            & 0,972 \\
Enfermedad cardíaca   & Medicina           &     303 & Bosque Aleatorio    & 0,839 \\
Titanic               & Histórico          &     891 & Bosque Aleatorio    & 0,977 \\
Renta adultos         & Socioeconómico     & 48.842  & LightGBM            & 0,750 \\
Marketing bancario    & Financiero         & 45.211  & LightGBM            & 0,749 \\
\bottomrule
\end{tabular}
\end{table}

\begin{table}[H]
\centering
\caption{Resultados del benchmark de regresión (5 \textit{datasets}).
 Métrica: RMSE sobre el conjunto de validación (5-fold).}
\label{tab:pool-regression}
\small
\begin{tabular}{llrll}
\toprule
\textbf{Dataset} & \textbf{Dominio} & \textbf{Filas} & \textbf{Modelo campeón} & \textbf{RMSE} \\
\midrule
Diabetes              & Medicina           &     442 & Ridge     & 55,38 \\
Eficiencia energética & Ingeniería         &     768 & Ridge     &  5,20 \\
Resistencia hormigón  & Construcción       &   1.030 & CatBoost  &  4,52 \\
Alquiler bicicletas   & Transporte         & 17.389  & CatBoost  & 38,25 \\
Vivienda California   & Inmobiliario       & 20.640  & LightGBM  &  0,45 \\
\bottomrule
\end{tabular}
\end{table}

Los resultados de clasificación reflejan el comportamiento esperado de las reglas
estáticas del planificador: los \textit{datasets} muy pequeños (Iris, Enfermedad
cardíaca) eligen modelos de baja complejidad (Regresión Logística, Bosque Aleatorio),
mientras que los conjuntos medianos con miles de instancias (Renta adultos, Marketing
bancario) favorecen LightGBM.
Las métricas inferiores en estos últimos (macro-F1 $\approx 0{,}75$) se deben al
desbalanceo de clases —la clase positiva representa el 11\,\% en Marketing bancario—
y no a un defecto del sistema de selección.
En regresión destaca que Ridge gana en los dos \textit{datasets} más pequeños
(Diabetes, Eficiencia energética), confirmando la regla de regularización para
conjuntos con pocas muestras.

% ============================================================
\subsection{Comportamiento agéntico: casos de estudio}
\label{sec:casos-estudio}
% ============================================================

Mientras que la Sección~\ref{sec:pool-evaluacion} evalúa el \textit{pool} de forma
agregada y determinista, esta sección examina el sistema \emph{en funcionamiento}:
seis casos de estudio a través del pipeline interactivo completo, cada uno diseñado para
validar un aspecto distinto del comportamiento agéntico. Todas las ejecuciones se
realizaron con los modelos \texttt{gpt-5.4-mini} (validador y planificador) y
\texttt{gpt-5.4-nano} (informe de auditoría) sobre la capa gratuita de GitHub~Models.
La Tabla~\ref{tab:casos-resumen} resume las cuatro ejecuciones de pronóstico; las
métricas de error son SMAPE (\textit{Symmetric Mean Absolute Percentage Error}),
adimensional y comparable entre escalas.

\begin{table}[H]
\centering
\caption{Resumen de los casos de estudio de pronóstico.
 Sim.\ = similitud de la experiencia recuperada más relevante;
 RMSE y SMAPE sobre el conjunto de prueba retenido.
 RF = \textit{Random Forest}; XGB = XGBoost.}
\label{tab:casos-resumen}
\small
\setlength{\tabcolsep}{4pt}
\begin{tabular}{lllrlll}
\toprule
\textbf{Caso} & \textbf{Dataset} & \textbf{Frec./Filas}
  & \textbf{Sim.} & \textbf{Experiencia top}
  & \textbf{Campeón} & \textbf{SMAPE} \\
\midrule
6.2.1 & Revenue mensual    & MS / 120 & 1,00 & Air Passengers      & AutoARIMA & 1,9\,\% \\
6.2.1 & Air Passengers     & MS / 144 & 0,83 & Air Passengers      & AutoARIMA & 4,0\,\% \\
6.2.2 & Grid Demand        & W / 156  & 0,82 & Air Passengers      & AutoARIMA & 2,4\,\% \\
6.2.3 & Bakery diario      & D / 1.000& 0,89 & Gold Macro (RF)     & ETS       & 3,4\,\% \\
\bottomrule
\end{tabular}
\end{table}

% ── 6.2.1  Recuperación del pool ────────────────────────────────────────────

\subsubsection{Recuperación de experiencias y razonamiento del planificador}
\label{sec:caso-retrieval}

El primer caso valida las afirmaciones~(ii) y~(iii): el \textit{pool} es consultado
realmente por el planificador y la función de similitud identifica priors relevantes.
Se ejecutaron dos \textit{datasets} mensuales distintos. El primero, un conjunto
sintético de ingresos mensuales (\texttt{small\_monthly\_revenue}, 120 observaciones),
produjo un perfil de \textit{dataset} idéntico en cubos al de Air~Passengers, de modo
que la herramienta \texttt{retrieve\_experiences} devolvió esa experiencia con
similitud~\textbf{1,00}. El segundo, el propio Air~Passengers cargado a través de la
interfaz, recuperó su experiencia homónima del \textit{pool} con similitud~0,83 —no~1,00
porque la ejecución interactiva incorpora un \textit{join} con contexto de aviación que
altera ligeramente el perfil.

En ambos casos el planificador citó la experiencia recuperada como evidencia primaria y
recomendó el mismo conjunto conservador de candidatos estadísticos —AutoARIMA, ETS,
\textit{seasonal\_naive} y \textit{naive}—, y el ejecutor seleccionó AutoARIMA como
campeón, coincidiendo con el modelo ganador del benchmark original
(Tabla~\ref{tab:pool-forecasting}). El análisis del planificador para el caso de
ingresos mensuales ilustra el patrón de razonamiento:

\begin{quote}\small\itshape
``El registro y las reglas favorecen modelos estadísticos bajo histórico corto, y la
experiencia de mayor relevancia también seleccionó AutoARIMA como mejor. Recomiendo por
tanto un conjunto conservador de candidatos estadísticos\ldots\ Rechazo los modelos
supervisados con \textit{lag} porque la regla de histórico corto los penaliza, aunque
algunas experiencias de similitud media favorecieron métodos de árboles en otros
\textit{datasets}.''
\end{quote}

Este fragmento evidencia el mecanismo central del sistema: la decisión no es una
alucinación del modelo, sino una síntesis trazable de tres fuentes —el registro de
modelos, las reglas estáticas y la experiencia empírica recuperada— que el planificador
pondera explícitamente, incluyendo el reconocimiento del conflicto con experiencias que
favorecieron otras familias.

% ── 6.2.2  Joins ────────────────────────────────────────────────────────────

\subsubsection{Descubrimiento de \textit{joins} multi-tabla}
\label{sec:caso-joins}

El caso de Grid~Demand valida la afirmación~(i) sobre la robustez del agente de
validación y, en particular, su capacidad de descubrimiento de \textit{joins}. Se
cargaron cuatro ficheros simultáneamente: la serie base de demanda eléctrica semanal
(\texttt{grid\_demand}, 156~filas, con la variable objetivo \texttt{demand\_mwh} y el
índice temporal) y tres tablas secundarias —registro de calendario, mezcla de
generación y datos meteorológicos—. Ninguna comparte el nombre de la columna clave: la
serie base usa \texttt{timestamp\_week}, mientras que las secundarias usan
\texttt{week\_start\_date}, \texttt{week\_ref} y \texttt{obs\_date} respectivamente.

El agente de validación, sin intervención humana, infirió los tres \textit{joins} de
tipo \textsc{left} uno-a-uno, los evaluó (cobertura del 100\,\%, multiplicador de filas
$1{,}00\times$, riesgo bajo) y los ejecutó, materializando un \textit{dataset} de
156~filas $\times$ 8~columnas con cinco variables exógenas añadidas
(Figura~\ref{fig:cap06-join-plan}). El plan de \textit{join} se presenta al revisor
humano en la puerta de aprobación con su nivel de confianza (``alta'') y la justificación
de cada candidato, materializando la afirmación~(iv): el humano \emph{confirma}, no
diseña.

\begin{figure}[H]
\centering
% Fuente: resultados/grid_demand/join_plan.png
\includegraphics[width=\textwidth]{figuras/cap06/grid_demand_join_plan.png}
\caption{Plan de \textit{join} inferido para Grid~Demand, tal como se presenta en la
 puerta de aprobación del \textit{dataset}. El agente descubre tres \textit{joins}
 uno-a-uno con cobertura total y riesgo bajo, alineando claves con nombres distintos
 (\texttt{timestamp\_week} $\leftrightarrow$ \texttt{week\_start\_date},
 \texttt{week\_ref}, \texttt{obs\_date}).}
\label{fig:cap06-join-plan}
\end{figure}

Pese a disponer de variables exógenas tras el \textit{join}, el planificador mantuvo un
plan conservador (AutoARIMA, ETS, \textit{seasonal\_naive}, \textit{naive}) porque el
histórico es corto (156~observaciones $<$ 200): el sistema prioriza la regla de
histórico corto sobre la disponibilidad de exógenas, y AutoARIMA resultó campeón.

% ── 6.2.3  Exógenas ─────────────────────────────────────────────────────────

\subsubsection{Adaptación de familia de modelos con variables exógenas}
\label{sec:caso-exog}

El caso de la panadería diaria (\texttt{medium\_daily\_bakery}, 1.000~observaciones)
valida la afirmación~(ii) en su forma más exigente: que la selección de modelos
\emph{cambia de familia} cuando la señal de los datos lo justifica. A diferencia de los
casos anteriores, este \textit{dataset} combina un histórico medio con variables
exógenas de calendario de valor futuro conocido (\texttt{is\_weekend},
\texttt{day\_of\_week}).

La recuperación del \textit{pool} devolvió como experiencias más similares conjuntos en
los que ganaron modelos supervisados —Gold~Macro (\textit{Random Forest}, sim.~0,89),
S\&P~500 (XGBoost, sim.~0,89) y Metro~Minneapolis (\textit{Random Forest}, sim.~0,82)—.
En consecuencia, y a diferencia de los casos puramente estadísticos, el planificador
\emph{amplió} el conjunto de candidatos para incluir \texttt{random\_forest\_forecaster}
y \texttt{xgboost\_forecaster} junto a ETS y AutoARIMA. Esto demuestra que el cambio de
familia no está codificado de forma rígida, sino que emerge de la combinación de la
disponibilidad de exógenas en el perfil y de los priors supervisados recuperados.

El resultado final es matizado y honesto: aunque el planificador ofreció candidatos
supervisados, el campeón empírico en validación fue ETS (RMSE~16,71; SMAPE~3,4\,\%),
que capturó la estacionalidad semanal del consumo de la panadería mejor que los modelos
de árboles en este horizonte. El sistema entrena todos los candidatos y deja que la
validación cruzada decida; el valor del planificador no es acertar siempre el campeón,
sino \emph{acotar} el espacio de búsqueda a candidatos plausibles.

% ── 6.2.4  Robustez del controlador ─────────────────────────────────────────

\subsubsection{Robustez del controlador ante entradas malformadas}
\label{sec:caso-robustez}

Este caso valida la afirmación~(i): el controlador determinista alcanza siempre un
estado terminal limpio e informativo, sin caídas ni corrupción silenciosa. Se
inyectaron tres tipos de fallo y se observó dónde los detecta el sistema
(Tabla~\ref{tab:robustez}).

\begin{table}[H]
\centering
\caption{Escenarios de error forzado y comportamiento observado del controlador.
 Todos alcanzan un estado terminal limpio con un mensaje informativo.}
\label{tab:robustez}
\small
\begin{tabular}{p{3.0cm}p{3.4cm}p{6.4cm}}
\toprule
\textbf{Fallo inyectado} & \textbf{Detectado en} & \textbf{Mensaje al usuario (verbatim)} \\
\midrule
Columna objetivo 100\,\% NaN
  & \texttt{data\_validator} (semántico)
  & \textit{``Data validation failed after auto-fix attempt''} \\
\addlinespace
Columna objetivo declarada pero ausente
  & \texttt{data\_validator} (contrato de esquema)
  & \textit{``target\_column `sales' not found in dataset''} \\
\addlinespace
15 filas para horizonte 12
  & \texttt{planner} (chequeo de capacidad)
  & \textit{``Forecasting capacity check failed: need $\geq$ 60 observations for horizon 12, have 15''} \\
\bottomrule
\end{tabular}
\end{table}

Los tres escenarios son detectados en \emph{etapas distintas} del pipeline, lo que
demuestra una defensa en profundidad: el contrato de esquema (puramente determinista)
captura la columna ausente antes incluso de invocar al agente; el agente de validación
captura el objetivo degenerado (todo~NaN) tras intentar la auto-reparación; y el
planificador captura la insuficiencia de histórico en su chequeo de capacidad, ya
superada la puerta de aprobación del \textit{dataset}. En los tres casos el controlador
enruta a \texttt{END} con \texttt{error\_message} establecido, y el pipeline emite un
evento \texttt{run\_complete} con el prefijo \textit{``Pipeline did not complete
end-to-end''} en lugar de un falso éxito. Ningún escenario produjo una excepción no
gestionada ni dejó artefactos corruptos en el \textit{pool}.

% ── 6.2.5  Generalidad ──────────────────────────────────────────────────────

\subsubsection{Generalidad del sistema}
\label{sec:caso-generalidad}

Aunque el foco de este capítulo es el pronóstico, la arquitectura agente–planificador–
ejecutor es agnóstica al dominio. La evidencia ya se presentó en la
Sección~\ref{sec:pool-clasif-regr}: ante \textit{datasets} tabulares no temporales, el
mismo planificador aplica reglas distintas —selecciona Regresión Logística o Bosque
Aleatorio para conjuntos pequeños y LightGBM para los de decenas de miles de filas— sin
proponer nunca un modelo de pronóstico. El cambio de tipo de problema se propaga desde
el esquema cargado por el usuario a través del estado compartido, y ni el controlador ni
el ejecutor requieren modificación. Las mejoras introducidas en el dominio de pronóstico
no degradan el comportamiento en clasificación o regresión.
El caso siguiente hace concreta esta generalidad con una ejecución de regresión que, a
diferencia de los \textit{datasets} tabulares limpios del \textit{benchmark}, ejercita
el camino completo de ingeniería de datos sobre fuentes multi-tabla con integridad
referencial imperfecta.

% ── 6.2.6  Joins imperfectos y características categóricas ───────────────────

\subsubsection{Joins imperfectos y características categóricas en regresión}
\label{sec:caso-joins-imperfectos}

Este último caso es el más exigente para el agente de validación porque concentra en una
sola subida los tres frentes que un escenario multi-tabla real plantea: integridad
referencial imperfecta, características categóricas que requieren imputación y
codificación, y columnas identificadoras que no deben convertirse en variables
predictoras. Es además un problema de regresión, por lo que materializa de forma concreta
la generalidad argumentada en la Sección~\ref{sec:caso-generalidad}.

Se cargaron cuatro ficheros: una tabla de hechos de ventas minoristas
(\texttt{retail\_sales}, 320~filas, objetivo \texttt{revenue\_eur}) y tres tablas de
dimensión —catálogo de productos, directorio de tiendas y segmentos de cliente— cuyas
claves tienen nombres distintos de los de la tabla de hechos
(\texttt{product\_id}$\,\leftrightarrow\,$\texttt{sku},
\texttt{store\_id}$\,\leftrightarrow\,$\texttt{store\_code},
\texttt{segment\_code}$\,\leftrightarrow\,$\texttt{segment\_code}). El descubrimiento de
\textit{joins} alineó las tres claves de forma semántica y reportó, para cada una, las
métricas que distinguen tres regímenes de integridad referencial muy distintos
(Tabla~\ref{tab:joins-imperfectos}).

\begin{table}[H]
\centering
\caption{\textit{Joins} descubiertos en \texttt{retail\_sales}.
 \textit{Cobertura izq.}\ = fracción de filas de la tabla de hechos con correspondencia;
 \textit{Contención} = relación entre el conjunto de claves foráneas (FK) y el de claves
 primarias (PK) de la dimensión.}
\label{tab:joins-imperfectos}
\small
\setlength{\tabcolsep}{4pt}
\begin{tabular}{lllll}
\toprule
\textbf{Dimensión} & \textbf{Clave (hecho\,$\leftrightarrow$\,dim.)}
  & \textbf{Cobertura izq.} & \textbf{Contención} & \textbf{Acción} \\
\midrule
Producto  & \texttt{product\_id}$\,\leftrightarrow\,$\texttt{sku}
          & 100\,\%   & FK\,$\subsetneq$\,PK (40/60) & \textit{Join} limpio \\
Tienda    & \texttt{store\_id}$\,\leftrightarrow\,$\texttt{store\_code}
          & 71,6\,\%  & parcial + huérfanas          & \textit{Join} + imputación + dedup. \\
Segmento  & \texttt{segment\_code}$\,\leftrightarrow\,$\texttt{segment\_code}
          & 100\,\%   & FK\,$\subsetneq$\,PK (4/6)    & \textit{Join} limpio \\
\bottomrule
\end{tabular}
\end{table}

La distinción que el sistema establece es clave. Las dimensiones de producto y segmento
presentan \emph{contención}: la dimensión contiene más claves de las que la tabla de
hechos utiliza (40 de 60~SKU; 4 de 6~segmentos), pero \emph{todas} las filas de hechos
encuentran correspondencia (cobertura izquierda del 100\,\%). Es el caso benigno
FK\,$\subsetneq$\,PK: un \textit{join} \textsc{left} limpio sin valores nulos. El
directorio de tiendas, en cambio, es el caso problemático: solo cubre el 71,6\,\% de las
filas (229 de 320), dejando 91~filas \emph{huérfanas} cuyos atributos de tienda
(\texttt{region}, \texttt{store\_size\_sqm}, \texttt{store\_type}) quedan nulos tras el
\textit{join}; además contiene una clave duplicada (\texttt{ST05} aparece dos veces), que
el descubrimiento señala como riesgo de cardinalidad y resuelve deduplicando la dimensión
para que el multiplicador de filas se mantenga en $1{,}00\times$ y la tabla de hechos no
se infle. Que el sistema separe estos dos regímenes —cobertura parcial benigna frente a
claves foráneas huérfanas— descansa en reportar por separado la cobertura izquierda y la
contención, y en aplicar un umbral de aviso (cobertura $<$ 80\,\%) que solo dispara la
alerta sobre el \textit{join} de tiendas.

El plan de \textit{join}, con su cobertura y su nivel de riesgo por candidato, se presenta
al revisor en la puerta de aprobación del \textit{dataset} (Figura~\ref{fig:cap06-join-imperfecto}),
de nuevo bajo el principio de que el humano \emph{confirma} la integridad imperfecta
detectada, no la diagnostica.

\begin{figure}[H]
\centering
% Fuente: captura de la tarjeta de aprobación del dataset (run retail_sales_partial_join)
\includegraphics[width=\textwidth]{figuras/cap06/retail_join_imperfecto.png}
\caption{Plan de \textit{join} para \texttt{retail\_sales} en la puerta de aprobación. El
 agente reporta la cobertura parcial del directorio de tiendas (71,6\,\%) y las filas
 huérfanas resultantes, junto a los dos \textit{joins} limpios de producto y segmento.}
\label{fig:cap06-join-imperfecto}
\end{figure}

Tras la aprobación, la cadena determinista resolvió los tres frentes sin intervención
adicional. Primero, el agente de validación \emph{imputó} las 91~filas huérfanas —los
atributos categóricos \texttt{region} y \texttt{store\_type} con la moda, y el numérico
\texttt{store\_size\_sqm} con la media— y revalidó el \textit{dataset} con éxito.
Segundo, el ejecutor \emph{descartó} la columna \texttt{txn\_id}: marcada como
\texttt{unique:\,true} en el esquema y no estructural (ni objetivo, ni índice temporal),
es un identificador puro que no debe filtrarse como ruido de alta cardinalidad al modelo.
Tercero, el \textit{pipeline} del modelo \emph{codificó} de forma ordinal las columnas
categóricas restantes (\texttt{category}, \texttt{brand}, \texttt{region},
\texttt{store\_type}, \texttt{segment\_name}) antes del ajuste, dejando intactas las
numéricas. El orden importa: la imputación de cadenas precede a la codificación, de modo
que el codificador nunca recibe nulos.

El campeón seleccionado en validación fue Ridge (RMSE~285,26; MAE~210,32; R$^2 \approx
0{,}50$ sobre el conjunto de prueba retenido). Es relevante que la puerta de evaluación
determinista \emph{retuviera} la promoción: con R$^2 = 0{,}50 < 0{,}70$ (umbral aplicado),
el sistema no ofreció la puerta de despliegue y alcanzó un estado terminal limpio sin
promover el modelo. Esto no es un fallo sino la conducta honesta esperada: el objetivo
sintético tiene una señal aprendible limitada por construcción, y el sistema declina
promover un modelo por debajo del umbral en lugar de enmascararlo. Las cuatro pestañas de
la interfaz —\textit{dataset}, planificador, modelo y auditoría— se completaron en verde,
documentando un recorrido íntegro pese a la decisión de no promover.

En conjunto, este caso ejercita toda la superficie de ingeniería de datos que una subida
multi-tabla real presenta —emparejamiento semántico de claves con nombres distintos, tres
regímenes de integridad referencial en una sola subida, cardinalidad con clave duplicada,
mezcla de características categóricas y numéricas, e higiene de identificadores— y la
columna vertebral determinista la resolvió en su totalidad sin más intervención humana que
la única puerta de confirmación. Es la evidencia más completa de la afirmación~(i): la
robustez del controlador y del agente de validación no se limita a entradas malformadas
(Sección~\ref{sec:caso-robustez}), sino que abarca la integridad referencial parcial,
el escenario más común y silencioso en datos tabulares del mundo real.

% ============================================================
\subsection{Coste agéntico y supervisión humana}
\label{sec:coste-agente}
% ============================================================

La inteligencia agéntica tiene un coste en tiempo y en \textit{tokens}. Esta sección lo
cuantifica a partir de la instrumentación del propio sistema (tarjeta de coste y
\textit{EventLog}) y argumenta que el compromiso es favorable.

% ── 6.3.1  Desglose ─────────────────────────────────────────────────────────

\subsubsection{Desglose de coste por nodo}
\label{sec:coste-desglose}

La Tabla~\ref{tab:coste-desglose} desglosa, para tres casos de pronóstico
representativos —Revenue, Grid y Bakery—, el
tiempo de cómputo por nodo (excluyendo las pausas de revisión humana en las puertas
HITL) y el coste monetario estimado. El tiempo se mide a partir de las marcas temporales
del \textit{EventLog}; el coste, a partir del recuento de \textit{tokens} y la tabla de
precios local (\texttt{gpt-5.4-mini}: \$0,75/\$4,50 por millón de \textit{tokens} de
entrada/salida; \texttt{gpt-5.4-nano}: \$0,20/\$1,25). Se omite Air~Passengers del
desglose porque no aporta un régimen de coste distinto: en su variante con \textit{join}
coincide con Grid y en su variante de fichero único con Revenue. Su aportación específica
—la validación de la recuperación del \textit{pool}— se documenta en
la Sección~\ref{sec:caso-retrieval}.

\begin{table}[H]
\centering
\caption{Desglose de tiempo de cómputo (segundos) y coste por ejecución. Los nodos
 \texttt{executor}, \texttt{evaluation} y \texttt{deployer} son deterministas; el resto
 son agénticos (LLM). El tiempo excluye las pausas humanas en las puertas HITL.
 Las tres columnas promedian 3 ejecuciones; los tiempos por nodo
 se reconstruyen a partir de las marcas de tiempo de los eventos del \textit{run}.}
\label{tab:coste-desglose}
\small
\setlength{\tabcolsep}{5pt}
\begin{tabular}{lrrrr}
\toprule
\textbf{Nodo} & \textbf{Tipo}
  & \textbf{Revenue} & \textbf{Grid} & \textbf{Bakery} \\
\midrule
\texttt{data\_validator} & LLM    & 23,6 & 44,2 & 16,9 \\
\texttt{planner}         & LLM    & 60,9 & 55,4 & 120,9 \\
\texttt{executor}        & Determ.& 17,3 & 20,7 & 95,9 \\
\texttt{evaluation}      & Determ.& 0,1 & 0,1 & 0,1 \\
\texttt{report\_writer}  & LLM    & 5,6 & 7,7 & 6,5 \\
\midrule
\textbf{Total cómputo}   &        & 107,5 & 128,1 & 240,1 \\
\textbf{Cuota LLM}       &        & 84\,\% & 84\,\% & 60\,\% \\
\midrule
\textbf{Coste (USD)}     &        & 0,091 & 0,175 & 0,137 \\
\bottomrule
\end{tabular}
\end{table}

Tres observaciones se desprenden de la tabla. Primero, el coste monetario por ejecución
completa es muy bajo —entre \$0,09 y \$0,18—, dominado por el agente de validación
(cuyo histórico de mensajes crece con cada ronda de herramientas) y por el planificador.
Segundo, la cuota del tiempo de cómputo atribuible a los nodos LLM oscila entre el 60\,\%
y el 84\,\%: un pipeline determinista tradicional solo pagaría la fracción restante
(\texttt{executor}\,+\,\texttt{evaluation}). Ese sobrecoste agéntico compra exactamente
las cuatro capacidades que justifican la arquitectura: selección adaptativa de modelos,
acumulación de experiencia, informes de auditoría en lenguaje natural y las puertas de
supervisión humana. Tercero, el tiempo del \texttt{executor} \emph{sí} escala con la
complejidad: 95,9~s en la panadería (1.000~filas y candidatos supervisados de árboles)
frente a 17--21~s en Revenue y Grid, de histórico corto y candidatos estadísticos.

\paragraph{Reintentos del agente como factor de inflación.}
Las cifras de un nodo agéntico pueden además inflarse por \emph{reintentos}: cuando la
salida estructurada de un agente no supera la validación determinista de su esquema, el
nodo se reejecuta dentro de su presupuesto de intentos. Es ilustrativo el caso de la
panadería, cuya columna promedia 3 ejecuciones: en 2 de ellas el planificador reintentó
—su salida citaba una evidencia de tipo \texttt{dataset\_profile} con \texttt{source\_id}
no nulo, rechazada por el validador— y ese reintento prácticamente \textbf{duplica tanto
el tiempo como el coste} de ese nodo, de $\sim$77~s y \$0,07 sin reintento a $\sim$140~s
y \$0,12 con él. El reintento es, por tanto, el principal factor de variabilidad de coste
entre ejecuciones de un mismo \textit{dataset}, por encima del tamaño o el tipo de
problema. Conviene subrayar que el reintento no es un fallo del sistema sino su garantía
de integridad funcionando: la validación determinista nunca acepta una salida malformada,
y prefiere pagar un segundo intento a propagar un plan inválido aguas abajo.

\paragraph{Nota honesta sobre la latencia.}
Los tiempos de los nodos LLM están acotados por la latencia de la API y los límites de
la capa gratuita de GitHub~Models (150 peticiones diarias por modelo), no por el cómputo
propio del sistema. En un nivel de pago, las fases agénticas serían considerablemente más
rápidas; las cifras de la Tabla~\ref{tab:coste-desglose} deben leerse como una cota
superior conservadora.

% ── 6.3.2  Independencia del tamaño ─────────────────────────────────────────

\subsubsection{Independencia del coste agéntico respecto al tamaño del dataset}
\label{sec:coste-tamano}

Una propiedad de diseño con consecuencias importantes para la escalabilidad es que el
coste en \textit{tokens} de los nodos agénticos es \emph{prácticamente independiente del
número de filas} del \textit{dataset}. La evidencia más limpia es el agente de
validación: procesó un \textit{dataset} de 120~filas con 36.557 \textit{tokens} de
entrada y uno de 1.000~filas con 32.360 —ocho veces más datos y, de hecho, ligeramente
\emph{menos} \textit{tokens}—.

La razón es estructural: los agentes \textbf{nunca reciben los datos en bruto}. El CSV
completo lo procesan herramientas deterministas (pandas) sobre disco, y al modelo solo
llegan resúmenes compactos. La salida íntegra que la herramienta \texttt{load\_dataset}
envió al LLM para el \textit{dataset} de 1.000~filas ocupó 451~caracteres:

\begin{lstlisting}[caption={Resumen que recibe el LLM para un dataset de 1.000
 filas: recuentos, nombres de columna, tipos y una muestra fija de tres filas.},
 label=lst:load-dataset, basicstyle=\footnotesize\ttfamily, breaklines=true]
{"row_count": 1000, "column_count": 4,
 "column_names": ["date","units_sold","is_weekend","day_of_week"],
 "dtypes": {"date":"object","units_sold":"int64", ...},
 "head": [ {3 filas de muestra} ]}
\end{lstlisting}

El recuento de filas es un único entero y la muestra (\texttt{head}) tiene un tamaño
fijo de tres filas, sea el \textit{dataset} de 120 o de 10~millones de filas. El
planificador, análogamente, no ve filas: opera sobre el \emph{perfil} discretizado del
\textit{dataset} (frecuencia, estacionalidad y tendencia detectadas, tramo de longitud de
histórico, etc.), también de tamaño constante. Es el principio ``determinista primero''
hecho coste: los datos pesados permanecen en Python y el LLM solo ve una abstracción
compacta e independiente del volumen.

En consecuencia, lo que sí mueve el recuento de \textit{tokens} no es el tamaño del
\textit{dataset} sino la \emph{complejidad agéntica}: el número de rondas de herramientas
—cada ronda reenvía el historial de mensajes acumulado—, el número de ficheros y
\textit{joins} (el agente de validación de Grid~Demand alcanzó 140.176 \textit{tokens} de
entrada por sus cuatro ficheros y diez rondas) y los reintentos del planificador (la
Sección~\ref{sec:coste-variabilidad} lo cuantifica). El único efecto del propio
\textit{dataset} es lineal y débil en el número de \emph{columnas} —que alarga las listas
de nombres y tipos—, nunca en el número de filas.

Combinada con la Sección~\ref{sec:coste-desglose}, esta propiedad permite caracterizar el
coste total del pipeline como la suma de dos términos de naturaleza opuesta: un
\textbf{sobrecoste agéntico aproximadamente constante} (validación, planificación y
auditoría, $\sim$150~s e independiente del tamaño) y un \textbf{coste de entrenamiento
determinista que escala con los datos} ($O(\text{filas}\times\text{particiones}\times
\text{ensayos}\times\text{modelos})$, como en cualquier sistema de AutoML). La inteligencia
agéntica, por tanto, añade un coste fijo que no se degrada con la escala; el cuello de
botella en \textit{datasets} grandes reside en el núcleo determinista de entrenamiento,
compartido por cualquier \textit{pipeline} de ML, y no en la capa de agentes —que incluso
lo alivia al podar el conjunto de candidatos—.

% ── 6.3.3  Variabilidad ─────────────────────────────────────────────────────

\subsubsection{Variabilidad del coste agéntico}
\label{sec:coste-variabilidad}

Una propiedad inesperada y relevante de los componentes agénticos es que su coste
\emph{no es determinista}, aun cuando la decisión sí lo es. La Tabla~\ref{tab:variabilidad}
recoge cuatro ejecuciones del mismo \textit{dataset} (la panadería) con el mismo modelo
y la misma configuración: la familia de candidatos y el campeón seleccionado se mantienen
estables entre repeticiones, pero la duración del planificador varía entre 84 y 131~s y
sus \textit{tokens} de salida entre 12,8k y 26,3k.

\begin{table}[H]
\centering
\caption{Cuatro repeticiones de la panadería: decisión estable, coste variable.
 \textit{Tokens} de salida del planificador (incluyen los de razonamiento).}
\label{tab:variabilidad}
\small
\begin{tabular}{lrrrl}
\toprule
\textbf{Ejecución} & \textbf{Validador (s)} & \textbf{Planner (s)}
  & \textbf{Tokens salida planner} & \textbf{Reintento} \\
\midrule
1 & 13,6 & 130,6 & 15,4k & no \\
2 & 4,1  & 84,3  & 14,5k & no \\
3 & 20,0 & 91,0  & 26,3k & sí \\
4 & 12,2 & 120,3 & 12,8k & no \\
\bottomrule
\end{tabular}
\end{table}

Esta variabilidad tiene dos fuentes. La primera es intrínseca a los modelos de
razonamiento: los \textit{tokens} de razonamiento privado (la cadena de pensamiento
oculta) constituyen más de la mitad de la salida del planificador y su volumen fluctúa
entre ejecuciones aunque la respuesta final sea equivalente; la inferencia en
\textit{hardware} compartido no es además bit-reproducible ni siquiera a temperatura
cero. La segunda fuente es discreta: si la salida del planificador no supera la
validación determinista de estructura (las comprobaciones \texttt{\_check\_*}), el nodo
reintenta con realimentación correctiva, duplicando aproximadamente el trabajo. La
ejecución~3 de la Tabla~\ref{tab:variabilidad} es precisamente un caso de reintento
(\texttt{retry\_ok}), lo que explica su mayor recuento de \textit{tokens}. La conclusión
para el diseño es importante: la \emph{decisión} del sistema es estable y reproducible
—gracias a la recuperación determinista del \textit{pool} y a las reglas—, mientras que
su \emph{coste} es estocástico y dominado por el razonamiento del planificador.

% ── 6.3.4  HITL ─────────────────────────────────────────────────────────────

\subsubsection{Supervisión humana en las puertas HITL}
\label{sec:coste-hitl}

El sistema interrumpe la ejecución en dos puertas de aprobación humana: tras la
validación del \textit{dataset} (incluido el plan de \textit{join}) y antes de promover
el modelo al registro. El argumento central de la afirmación~(iv) es que el coste de
esta supervisión escala \emph{inversamente} con la calidad de las decisiones del agente:
cuando el plan es correcto, la tarea del humano es confirmar —una revisión de segundos—,
no diseñar.

Las marcas temporales del \textit{EventLog} respaldan este argumento. En la puerta de
aprobación del \textit{dataset}, la pausa humana observada osciló entre los 20~s de la
panadería —donde el plan de validación era trivialmente correcto— y varios minutos en
los casos con \textit{join} (Air~Passengers, Grid~Demand), donde el revisor inspeccionó
el plan de \textit{join} inferido antes de aprobarlo. Esta diferencia es la evidencia
empírica del mecanismo: el humano dedica más tiempo precisamente cuando hay una decisión
no trivial que verificar (un \textit{join} multi-tabla), y casi ninguno cuando el agente
no ha tomado decisiones discutibles.

De aquí se sigue la principal limitación honesta del enfoque y su conexión con el
trabajo futuro. Para \textit{datasets} genuinamente distintos de todo lo presente en el
\textit{pool} —dominios nuevos, frecuencias inusuales, estructuras exógenas no vistas—,
el plan del planificador es menos fiable y el humano debe invertir más tiempo en la
puerta. El incentivo, por tanto, es seguir poblando el \textit{pool} de experiencias:
más priors implican planes mejores, puertas más ligeras y, en el límite, una operación
progresivamente más automática para los perfiles de \textit{dataset} comunes.
